\documentclass[10pt]{beamer}
\usetheme{metropolis}
\usepackage{graphicx}
\usepackage{listings}
\usepackage{booktabs}
\usepackage{xcolor}
\definecolor{codebg}{HTML}{F5F5F5}
\definecolor{codegreen}{HTML}{2E7D32}
\definecolor{codegray}{HTML}{757575}
\definecolor{codeblue}{HTML}{1565C0}
\lstdefinestyle{rstyle}{language=R,backgroundcolor=\color{codebg},basicstyle=\ttfamily\scriptsize,keywordstyle=\color{codeblue}\bfseries,commentstyle=\color{codegray}\itshape,breaklines=true,frame=single,rulecolor=\color{codegray!30},showstringspaces=false,morekeywords={library,sf,tmap,tm_shape,tm_polygons,tm_fill,tm_borders,tm_layout,ggplot,geom_sf,scale_fill_viridis,theme_void,labs,tmap_mode}}
\lstdefinestyle{pythonstyle}{language=Python,backgroundcolor=\color{codebg},basicstyle=\ttfamily\scriptsize,keywordstyle=\color{codeblue}\bfseries,commentstyle=\color{codegray}\itshape,breaklines=true,frame=single,rulecolor=\color{codegray!30},showstringspaces=false,morekeywords={import,from,as,gpd,plt,folium,Choropleth,plot,Map}}
\title{Choropleth Maps}
\subtitle{Workshop 5 --- Module 6}
\author{Prof.Asoc.\ Endri Raco}
\institute{Data Visualization for Data Scientists \\ UNYT Tirana}
\date{2025/26}
\begin{document}
\begin{frame}\titlepage\end{frame}
\begin{frame}{Module Roadmap}
\begin{enumerate}
\item Choropleth concept: regions coloured by data
\item Classification methods: equal interval, quantile, Jenks
\item Palette selection: sequential, diverging, qualitative
\item Implementation in R and Python
\item Five common pitfalls
\end{enumerate}
\begin{exampleblock}{Learning Outcome}
You will produce choropleth maps with appropriate classification and palette, using both static and interactive tools.
\end{exampleblock}
\end{frame}
\section{Concept}
\begin{frame}{Choropleth: Colour = Value per Region}
\begin{center}\includegraphics[width=0.92\textwidth]{figures_w05m06/choropleth_concept}\end{center}
\end{frame}
\begin{frame}{Palette Selection}
\begin{center}\includegraphics[width=0.88\textwidth]{figures_w05m06/palette_choice}\end{center}
\begin{alertblock}{Pro-Tip}
Sequential data $\rightarrow$ sequential palette. Deviation $\rightarrow$ diverging centred at 0. Categories $\rightarrow$ qualitative. Never qualitative for numeric data.
\end{alertblock}
\end{frame}
\section{Code}
\begin{frame}[fragile]{Choropleth in R (ggplot2 + sf, tmap)}
\begin{lstlisting}[style=rstyle]
library(sf); library(tmap)
districts <- st_read("albania.geojson")
# ggplot2 static
ggplot(districts) +
  geom_sf(aes(fill = population), color = "white") +
  scale_fill_viridis_c(name = "Pop.") +
  theme_void() + labs(title = "Albania: Population")
# tmap static
tm_shape(districts) +
  tm_fill("gdp", style = "jenks", n = 5, palette = "Blues") +
  tm_borders(col = "white") +
  tm_layout(title = "GDP per Capita")
# tmap interactive (one command switch!)
tmap_mode("view")
tm_shape(districts) + tm_fill("gdp", style = "jenks")
\end{lstlisting}
\end{frame}
\begin{frame}[fragile]{Choropleth in Python (geopandas, folium)}
\begin{lstlisting}[style=pythonstyle]
import geopandas as gpd
districts = gpd.read_file("albania.geojson")
# Static
fig, ax = plt.subplots(figsize=(6, 8))
districts.plot(column="population", cmap="plasma",
    legend=True, edgecolor="white", ax=ax,
    scheme="NaturalBreaks", k=5)
ax.axis("off")
# Interactive (folium)
import folium
m = folium.Map(location=[41.33, 19.82], zoom_start=7)
folium.Choropleth(
    geo_data="albania.geojson",
    data=districts,
    columns=["id", "population"],
    key_on="feature.properties.id",
    fill_color="YlOrRd").add_to(m)
m.save("choropleth.html")
\end{lstlisting}
\end{frame}
\section{Pitfalls}
\begin{frame}{Five Choropleth Pitfalls}
\begin{center}\includegraphics[width=0.82\textwidth]{figures_w05m06/pitfalls}\end{center}
\end{frame}
\section{Synthesis}
\begin{frame}{Key Takeaways}
\begin{enumerate}
\item Choropleths colour \textbf{regions} by a variable. Best for rates; misleading for raw counts (large regions dominate).
\item \textbf{Jenks} (natural breaks) is the default classification. Quantile for skewed data. 5--7 classes maximum.
\item \textbf{Sequential} palette for magnitude; \textbf{diverging} for deviation; never qualitative for numeric data.
\item \textbf{R}: \texttt{geom\_sf()} for static, \texttt{tmap\_mode("view")} for interactive.
\item \textbf{Python}: \texttt{geopandas.plot(scheme=)} for static, \texttt{folium.Choropleth()} for interactive.
\end{enumerate}
\end{frame}
\begin{frame}{Essential Readings}
\begin{itemize}
\item Lovelace, R. et al. (2019). \emph{Geocomputation with R}, Chapter 8.
\item Brewer, C.~A. (2005). \emph{Designing Better Maps}.
\item \texttt{colorbrewer2.org} --- palette tool for cartography.
\end{itemize}
\vspace{6pt}
\textbf{Next Module}: Point Maps \& Spatial Patterns (W05-M07)
\end{frame}
\end{document}
